\chapter{Metodologia}

Este capitulo describe el pipeline experimental completo usado en la etapa B2. La idea central fue mantener un protocolo reproducible y comparable entre familias de metodos, sin ajustar cada metodo de forma independiente para maximizar solo su mejor caso.

\section{Construccion de grafos}

Se evaluaron constructores de grafo de tipo vecindad, distancia y variantes adaptativas. En B2 se mantuvieron los metodos que mostraron estabilidad numerica en corrida extendida y que permitieron consolidar resultados sin interrupciones.

\section{Simulacion de canales faltantes}

Se trabajaron escenarios por region de electrodos (por ejemplo frontal, occipital y temporal lateral) con niveles de perdida de 10\%, 20\%, 30\% y 40\%. Este diseno intenta aproximar situaciones mas realistas que una mascara completamente aleatoria.

\section{Interpolacion y reconstruccion}

Las comparaciones se organizaron por familias:
\begin{itemize}
\item Baselines instantaneos/geometricos.
\item Metodos de suavidad en grafo (incluyendo variantes tipo GSP/Tikhonov).
\item Metodos TV/tiempo (incluyendo TRSS y variantes temporales).
\end{itemize}

Para la evaluacion se usaron MAE, RMSE, DTW y SNR con agregacion estadistica (media, desviacion estandar e intervalo de confianza aproximado) sobre corridas repetidas.
